%! Diagonalizing a Matrix — Unit 6, Lesson 6.2 — Group Activity
\documentclass[10pt]{article}
\usepackage{linalg-article}
\usepackage{linalg-boxes}
\newcommand{\TallMath}[1]{$\displaystyle #1\rule[-1.4em]{0pt}{3.2em}$}
\newcommand{\xx}{\mathbf{x}}
\newcommand{\zero}{\mathbf{0}}

\begin{document}

\pageheader{Unit 6, Lesson 6.2}{Group Activity}
\namepartnerperiod

\vspace{0.10in}

\begin{headlinebox}{blush}
\small\textbf{Diagonalizing a Matrix.} Line the eigenvectors into the columns of $X$ and the eigenvalues
down the diagonal of $\Lambda$. Then $AX=X\Lambda$, and (when $X$ is invertible)
$A=X\Lambda X^{-1}$. The payoff: $A^k=X\Lambda^k X^{-1}$ --- powers become powers of the $\lambda$'s. Show
every step with your partner.
\end{headlinebox}

\vspace{0.08in}

\begin{tcolorbox}[colback=white, colframe=black!40, title=\textbf{Tier R --- Remediate},
    fonttitle=\bfseries, arc=1.5mm, left=3mm, right=3mm, top=2mm, bottom=2mm, breakable]
Build the pieces, then assemble. A matrix $A$ has eigenvector $\begin{bsmallmatrix}1\\0\end{bsmallmatrix}$
with $\lambda=6$ and eigenvector $\begin{bsmallmatrix}1\\1\end{bsmallmatrix}$ with $\lambda=2$.
\begin{enumerate}[leftmargin=1.7em, itemsep=9pt]
  \item \textbf{Line them up.} Write $X$ (eigenvectors in the columns) and $\Lambda$ (eigenvalues on the diagonal).
        \[ X=\begin{bmatrix}\rule{0pt}{1.1em}\blank{0.6cm}&\blank{0.6cm}\\[3pt]\blank{0.6cm}&\blank{0.6cm}\end{bmatrix},\qquad
           \Lambda=\begin{bmatrix}\rule{0pt}{1.1em}\blank{0.6cm}&\blank{0.6cm}\\[3pt]\blank{0.6cm}&\blank{0.6cm}\end{bmatrix}. \]
  \item \textbf{Invert $X$} (Unit~2 formula). $X^{-1}=\blank{2.6cm}$.
  \item \textbf{Assemble $A=X\Lambda X^{-1}$.} Multiply $X\Lambda$ first (scales the columns), then times $X^{-1}$. \writelines{2}
\end{enumerate}
\end{tcolorbox}

\begin{tcolorbox}[colback=white, colframe=black!40, title=\textbf{Tier A --- Approaching Proficiency},
    fonttitle=\bfseries, arc=1.5mm, left=3mm, right=3mm, top=2mm, bottom=2mm, breakable]
Run the full method, then use it for a power. Let $A=\begin{bsmallmatrix} 4 & -2\\ 1 & 1 \end{bsmallmatrix}$.
\begin{enumerate}[leftmargin=1.7em, itemsep=9pt]
  \item \textbf{Diagonalize.}
        \begin{enumerate}[label=(\alph*), leftmargin=1.7em, itemsep=3pt]
          \item Find both eigenvalues from $\det(A-\lambda I)=0$, then an eigenvector for each.
          \item Write $X$, $\Lambda$, and $X^{-1}$, and confirm $A=X\Lambda X^{-1}$.
        \end{enumerate}
        \writelines{3}
  \item \textbf{Use it for a power.} Write $A^2=X\Lambda^2 X^{-1}$ and multiply it out. Check by squaring
        $A$ directly. \writelines{2}
\end{enumerate}
\end{tcolorbox}

\begin{tcolorbox}[colback=white, colframe=black!40, title=\textbf{Tier E --- Extension},
    fonttitle=\bfseries, arc=1.5mm, left=3mm, right=3mm, top=2mm, bottom=2mm, breakable]
\textbf{When diagonalization \emph{fails}.} Consider the shear $A=\begin{bsmallmatrix} 2 & 1\\ 0 & 2 \end{bsmallmatrix}$.
\begin{enumerate}[leftmargin=1.7em, itemsep=9pt]
  \item Find its eigenvalues from $\det(A-\lambda I)=0$. What do you notice? \writeline
  \item Solve $(A-2I)\xx=\zero$. How many \emph{independent} eigenvector directions are there? \writeline
  \item Explain why you \emph{cannot} build an invertible $X$ here --- so $A$ is \emph{not diagonalizable}.
        (Contrast: a matrix with two \emph{different} eigenvalues always can be.) \writelines{2}
\end{enumerate}
\end{tcolorbox}

\end{document}
